\section{Worum es geht}

Aus den Merkmalen eines Kreditnehmers soll vorhergesagt werden, ob er
ausfaellt. Das ist ein Klassifikationsproblem mit \emph{ungleichen Kosten}:
Ein verpasster Ausfall ist teuer, ein Fehlalarm kostet nur eine Pruefung.
Der Bericht misst daher nicht nur die Trennschaerfe, sondern leitet aus einem
Kostenmodell ab, wo die Entscheidungsschwelle liegen sollte.

Drei Fragen stehen im Vordergrund:

\begin{enumerate}
  \item Wie gut lassen sich Kreditnehmer ueberhaupt nach Risiko ordnen?
  \item Welche Merkmale haengen mit einem Ausfall zusammen -- und welche
        nur scheinbar?
  \item Wo sollte die Entscheidungsschwelle liegen, wenn ein verpasster
        Ausfall ein Vielfaches eines Fehlalarms kostet?
\end{enumerate}
